\scanpage{123}

$\Phi_n(T)$ polynôme minimal sur $\unsure{\mathbb Q|Q}$ d'une racine primitive $n$-ième $\zeta$ de 1

\margin{$n\in\mathbb N^*$}

$F_n(T)$ polynôme minimal de $\zeta+\zeta^{-1}$

Si $n\ne1,2$, on a $\Phi_n(T)=T^{\varphi'(n)}F_n(T+T^{-1})$ \quad $(\varphi'(n)=\varphi(n)/2)$

\margin{$\Phi_n,F_n\in\mathbb Z[T]$ pol. unitaires}

\[
\begin{aligned}
\deg\Phi_n&=\varphi(n)\quad\text{(indicateur d'Euler, }=\operatorname{card}(\mathbb Z/n\mathbb Z)^*\text{)}\\
\deg F_n&=\begin{cases}1&\text{si }n=1,2\quad(F_1(T)=T-2,\ F_2(T)=T+2)\\
\varphi(n)/2&\text{si }n\geq3\end{cases}
\end{aligned}
\]

\[
\prod_{\substack{d\in\mathbb N^*\\d\mid n}}\Phi_n(T)=T^n-1
\]

donne mode de calcul par récurrence \unsure{multiplicative} sur $n$.

\diagram{$\left\{\begin{array}{l}
F_1(\deleted{T}U)=\deleted{T}U-2\\
F_2(\deleted{T}U)=\deleted{T}U+2\\
F_3(\deleted{T}U)=\deleted{T}U+1\\
F_4(\deleted{T}U)=\deleted{T}U\deleted{\unsure{{}^2-2}}\\
F_5(\deleted{T}U)=\deleted{T}U^2+\deleted{T}U-1\\
F_6(\deleted{T}U)=\deleted{T}U-1\\
F_7(\deleted{T}U)=\deleted{T}U^3+\deleted{T}U^2-2\deleted{T}U-1\\
F_8(\deleted{T}U)=\deleted{T}U^2-2\\
F_9(\deleted{T}U)=\deleted{T}U^3-3\deleted{T}U+1\\
F_{10}(\deleted{T}U)=\deleted{T}U^2-\deleted{T}U-1\\
\ldots
\end{array}\right.\quad
\left\{\begin{array}{l}
\Phi_1(T)=T-1\\
\Phi_2(T)=T+1\\
\Phi_3(T)=T^2+T+1\\
\Phi_4(T)=T^2+1\\
\Phi_5(T)=T^4+T^3+T^2+T+1\\
\Phi_6(T)=T^2-T+1\\
\Phi_7(T)=T^6+T^5+T^4+T^3+T^{\unsure{1|2}}+T^{\unsure{2|1|{}}}+1\\
\Phi_8(T)=T^4+1\\
\Phi_9(T)=T^6+T^3+1\\
\Phi_{10}(T)=T^4-T^3+T^2-T+1\\
\ldots
\end{array}\right.$}

(donner aussi les cas $n=12,14,18\ldots$)

\[
\left\{\begin{aligned}
&\Phi_p(T)=1+T+\ldots+T^{p-1}\\
&\text{si }p\text{ premier}\\
&F_p(T)=S_0(T)+S_1(T)+\ldots+S_{p'}(T)
\end{aligned}\right.
\]

où $p=2p'+1$ i.e. $p-1=2p'$ ($p$ \unsure{nb|est} premier impair), et les

\[
S_{\unsure{n|p}}(\deleted{T}U)\in\mathbb Z[\deleted{T}U]\qquad(n\in\mathbb N^*)
\]

sont caractérisés par

\[
\left\{\begin{aligned}
&S_0(\deleted{T}U)=1\\
&T^n+T^{-n}=S_n\unsure{'}(T+T^{-1})\qquad n\geq1
\end{aligned}\right.
\]

donc

\[
\left\{\begin{aligned}
S_0(\deleted{T}U)&=1\\
S_1(\deleted{T}U)&=\deleted{T}U\\
S_2(\deleted{T}U)&=\deleted{T}U^2-2\\
S_3(\deleted{T}U)&=U^3-3U\\
&\ldots
\end{aligned}\right.
\]

Calcul récurrent des $S_i$ par formule

\[
\deleted{\unsure{S_n(T|S_n(T)}}\quad U^n=\sum_{0\leq i\leq[\frac n2]}\binom ni S_{n-2i}(U)\qquad n\geq1
\]

d'où

\[
S_n(U)=U^n-\sum_{1\leq i\leq[\frac n2]}\binom ni S_{n-2i}(U)
\]

% archive sheet number: 122
